\section*{MỞ ĐẦU}
\phantomsection\addcontentsline{toc}{section}{\numberline{}MỞ ĐẦU}
Số hóa quản trị nhân sự đã làm cho các hệ thống ghi nhận giờ vào, giờ ra và lịch làm việc trở nên phổ biến. Tuy nhiên, dữ liệu chấm công tại cửa hoặc dữ liệu đăng nhập máy tính chỉ phản ánh một phần trạng thái làm việc. Một nhân viên có thể chấm công rồi rời khu vực trong thời gian dài, ngược lại, trong một số môi trường kỹ thuật, nhân viên có thể làm việc trong vùng được phân công nhưng không thao tác liên tục trên máy tính. Vì vậy, việc ước lượng thời gian hiện diện tại khu vực làm việc là một nguồn thông tin bổ sung có giá trị cho công tác vận hành, miễn là khái niệm này được định nghĩa đúng và không bị đồng nhất với năng suất lao động.

Camera CCTV có ưu điểm là bao quát được nhiều người trong cùng một khung hình và có thể sử dụng hạ tầng sẵn có. Tuy nhiên, chuyển video thành dữ liệu thời gian không phải là bài toán phát hiện đối tượng đơn thuần. Hệ thống phải giải quyết đồng thời các vấn đề: phát hiện nhiều người, duy trì quỹ đạo ổn định khi che khuất, ánh xạ từng quỹ đạo với khu vực làm việc, hỗ trợ xác minh danh tính, loại bỏ dao động ở cấp khung hình, ghép các quan sát thành khoảng thời gian, và tránh tính trùng hoặc bỏ sót khi camera hoặc mạng tạm thời gặp lỗi. Các thách thức này khiến việc chỉ kiểm tra ``có người/không có người'' trên từng khung hình là không đủ.

Đồ án lựa chọn hướng \textit{theo dõi dựa trên phát hiện đối tượng}. YOLOv8 được dùng làm bộ phát hiện nền vì có thể phát hiện đồng thời người và điện thoại từ cùng một lần suy luận, ByteTrack được dùng để duy trì quỹ đạo qua chuỗi ảnh. ByteTrack tận dụng cả các quan sát phát hiện có độ tin cậy thấp trong bước liên kết nhằm phục hồi đối tượng bị che khuất ngắn hạn, thay vì loại bỏ toàn bộ chúng ngay từ đầu \cite{bytetrack}. Tư tưởng này phù hợp với cảnh CCTV văn phòng, nơi người có thể bị bàn ghế hoặc người khác che một phần.

Bên cạnh hiện diện, đồ án bổ sung theo dõi thời gian sử dụng điện thoại. Điện thoại chỉ được gán cho một nhân viên khi quan sát phát hiện điện thoại có quan hệ không gian đủ rõ với hộp bao/ROI của quỹ đạo người, sau đó một máy trạng thái theo thời gian xác nhận bắt đầu và kết thúc phiên sử dụng. Chính sách cho phép cấu hình một mức thời gian điện thoại trong ca. Chỉ phần vượt mức cho phép mới được trừ khỏi thời gian hiện diện để tạo ra \textit{thời gian làm việc hiệu lực}. Cách làm này bảo toàn dữ liệu gốc: hệ thống vẫn lưu riêng tổng hiện diện, tổng thời gian điện thoại, phần vượt mức và thời gian hiệu lực, nhờ đó mọi phép điều chỉnh đều có thể truy vết.

Một nguyên tắc xuyên suốt báo cáo là hệ thống không suy luận ``đang làm việc hiệu quả'' từ hình ảnh. Hình ảnh là dữ liệu cá nhân và việc nhận diện khuôn mặt là xử lý dữ liệu nhạy cảm, bởi vậy giải pháp cần giới hạn mục đích, quyền truy cập, thời hạn lưu trữ và khả năng hậu kiểm. Tại Việt Nam, Luật Bảo vệ dữ liệu cá nhân số 91/2025/QH15 có hiệu lực từ ngày 01/01/2026, do đó thiết kế hệ thống phải đặt yêu cầu bảo vệ dữ liệu ngay từ kiến trúc thay vì xem đây là chức năng bổ sung \cite{vnprivacy}.

Báo cáo trình bày cơ sở lý thuyết, khảo sát yêu cầu, thiết kế hệ thống, phương án cài đặt và quy trình đánh giá. Chương 6 sử dụng dữ liệu tổng hợp có thể tái tạo để kiểm tra phép tính chỉ số và các ràng buộc thời lượng. Kết quả tính toán được phân biệt với thực nghiệm trên camera, vốn cần video có nhãn tham chiếu độc lập và thông tin cấu hình đầy đủ.